\documentclass[11pt]{article}
\usepackage{amsmath, amssymb, amsthm}
\usepackage{geometry}
\usepackage{hyperref}
\usepackage{booktabs}

\geometry{margin=1.1in}

\newcommand{\code}[1]{\texttt{#1}}
\newcommand{\prox}{\operatorname{prox}}

\title{Two Algorithms, Side by Side}
\author{Optimal Transport / Schr\"odinger Bridge Project\\
Companion to \code{nonuniform\_grid\_xy\_variables.tex}, \code{nonuniform\_grid\_norms.tex}}
\date{}

\begin{document}
\maketitle

\begin{abstract}
Just two listings, kept deliberately apart: (I) the LADMM iteration as it actually runs in
\code{ladmm.py}, in pure math notation; (II) Goldstein--Li--Yuan's PDHG (Algorithm 1, the
plain iteration, and Algorithm 2, the adaptive step-size version), transcribed with every
symbol that would otherwise collide with a Part~I symbol renamed to something new -- so a
letter never means two different things across the two parts. No attempt yet to relate the
two, substitute one into the other, or say which symbol plays which role -- that's later,
once both sides are on the table and agreed.
\end{abstract}

\section*{Part I -- my LADMM, as coded}
Iterates $x^k,y^k,\delta^k$; scalars $\gamma_k>0$ (penalty) and $\tau_k>0$ (proximal
step) -- iteration-dependent, since the goal is to let them adapt (not yet wired up, but
the iteration below is written for that from the start rather than as a fixed-step
special case); linear/affine map $A$ (boundary conditions baked in); convex functions
$f_1,f_2$. One step, $k\to k+1$, using $\gamma_k,\tau_k$ throughout (not $\gamma_{k+1}$
or $\tau_{k+1}$ -- those aren't chosen until after this step, from its own residuals):
\begin{align}
r^k &:= Ax^k - y^k \notag\\
x^{k+1} &:= \prox_{(1/\tau_k)f_1}\Big(x^k - \tfrac1{\tau_k} A^\top(\gamma_k r^k - \delta^k)\Big) \notag\\
y^{k+1} &:= \prox_{f_2/\gamma_k}\Big(Ax^{k+1} - \delta^k/\gamma_k\Big) \notag\\
\delta^{k+1} &:= \delta^k - \gamma_k\big(Ax^{k+1}-y^{k+1}\big) \notag
\end{align}
where $\prox_{th}(u):=\arg\min_p\, h(p)+\tfrac1{2t}\|p-u\|^2$. (Constant $\gamma,\tau$ is
just $\gamma_k\equiv\gamma,\ \tau_k\equiv\tau$ -- today's actual behavior, and every
formula below specializes to what was sent before when that holds.)

\emph{The $x$-update's step, written out.} In general $\prox_{(1/\tau_k)f_1}$ genuinely
depends on $\tau_k$ -- this is the standard proximal-gradient step (step size $1/\tau_k$
on both the gradient move \emph{and} the prox), and everywhere below is written to hold
for that general case, since $f_1$ needn't be an indicator for a generic LADMM problem.
In \emph{this} project's own problem $f_1$ \emph{is} the FP-feasibility indicator, so
$\prox_{(1/\tau_k)f_1}(v)$ collapses to the plain projection onto the feasible set for
\emph{any} $\tau_k$ (indicators' proxes don't depend on the step) -- which is why earlier
passes of this note wrote bare $\prox_{f_1}$ without losing anything \emph{for this
specific problem}. Kept general here regardless, so nothing below relies on that
coincidence.

\emph{One simplification, flagged:} \code{ladmm.py} also has an over-relaxation
parameter $\alpha$ (in general, $\hat z^{k+1}:=\alpha\, Ax^{k+1}+(1-\alpha)y^k$ replaces
$Ax^{k+1}$ in the $y$- and $\delta$-updates above). Every run in this project uses
$\alpha=1$, which is exactly the $\hat z^{k+1}=Ax^{k+1}$ substitution already made above
-- so this is the iteration actually being run, not a special case chosen for convenience.

\section*{Part II -- Goldstein--Li--Yuan's PDHG, renamed to avoid Part I's letters}
Renaming rule: keep their symbol unless it already means something else in Part~I, in
which case pick an unused letter. $A,f_1,f_2^*$ turned out (Part III) to be the literal
same objects as their operator and functions, not just analogues, so they're written
directly below rather than kept as separate placeholder symbols -- only $u,v$ (their
$x,y$) need to stay distinct from Part~I's own $x,y$, which mean something else. Table
first, algorithm after.
\begin{center}
\begin{tabular}{@{}lll@{}}
\toprule
Their symbol & Clashes with Part I? & Used here instead\\
\midrule
$x$ (primal iterate) & yes -- Part I's $x$ & $u$\\
$y$ (dual iterate) & yes -- Part I's $y$ & $v$\\
$A$ (coupling operator) & no -- same operator as Part I's & $A$\\
$F,G$ (convex functions) & no -- $F=f_1$, $G=f_2^*$ exactly & $f_1,\ f_2^*$\\
$\tau_k$ (primal step) & yes -- Part I's $\tau$ & $\theta_k$\\
$\sigma_k$ (dual step) & no & $\sigma_k$ (unchanged)\\
$\alpha_k$ (adaptivity level) & yes -- Part I's $\alpha$ & $\omega_k$\\
$\Delta,\eta,s$ & no & unchanged\\
$p_k,d_k$ (residuals) & no ($D^k,P^k$ in \code{ladmm.py} are uppercase) & unchanged\\
\bottomrule
\end{tabular}
\end{center}

\paragraph{Algorithm 1 (plain PDHG).} Saddle point $\min_u\max_v f_1(u)+v^\top Au-f_2^*(v)$:
\begin{align}
u^{k+1} &= \prox_{\theta_k f_1}\big(u^k - \theta_k A^\top v^k\big) \notag\\
v^{k+1} &= \prox_{\sigma_k f_2^*}\big(v^k + \sigma_k A(2u^{k+1}-u^k)\big) \notag
\end{align}

\paragraph{Algorithm 2 (adaptive PDHG).} After each Algorithm-1 step, form
\begin{align}
p_{k+1} &= \Big|\tfrac1{\theta_k}(u^k-u^{k+1}) - A^\top(v^k-v^{k+1})\Big| \notag\\
d_{k+1} &= \Big|\tfrac1{\sigma_k}(v^k-v^{k+1}) - A(u^k-u^{k+1})\Big| \notag
\end{align}
then
\[
p_{k+1}>s\Delta\, d_{k+1}\ \Rightarrow\
\theta_{k+1}=\tfrac{\theta_k}{1-\omega_k},\ \sigma_{k+1}=\sigma_k(1-\omega_k),\ \omega_{k+1}=\eta\,\omega_k;
\]
\[
p_{k+1}<\tfrac{s}{\Delta}d_{k+1}\ \Rightarrow\
\theta_{k+1}=\theta_k(1-\omega_k),\ \sigma_{k+1}=\tfrac{\sigma_k}{1-\omega_k},\ \omega_{k+1}=\eta\,\omega_k;
\]
otherwise $\theta_{k+1}=\theta_k,\ \sigma_{k+1}=\sigma_k,\ \omega_{k+1}=\omega_k$. Required:
$\sigma_0\theta_0<\rho(A^\top A)^{-1}$. Their recommended defaults: $\omega_0=0.5$,
$\Delta=1.5$, $\eta=0.95$, $\theta_0=\sigma_0=1/\sqrt{\rho(A^\top A)}$.

\section*{Part III -- relating them}
The previous pass got stuck because Part II's Algorithm 1 extrapolates the wrong
variable to ever line up with Part I term-by-term. Fix suggested: reorder Algorithm 1
so the extrapolation is its own line, one step behind, built from a new variable $w$.
That resolves it completely -- not just ``closely related,'' an exact match.

Define, directly off Part I's own $x$-update bracket -- no elimination, this is
literally \code{ladmm.py}'s \texttt{gamma*r - delta}, already computed every iteration:
\[
w^k := \gamma_k r^k - \delta^k, \qquad u^k:=x^k,\qquad v^k:=-\delta^k.
\]

\subsection*{Algorithm 1, reordered}
\begin{align}
u^{k+1} &= \prox_{\theta_k f_1}\big(u^k-\theta_k A^\top w^k\big) \label{eq:u-reord}\\
v^{k+1} &= \prox_{\sigma_k f_2^*}\big(v^k+\sigma_k Au^{k+1}\big) \label{eq:v-reord}\\
w^{k+1} &= (1+\rho_{k+1})\,v^{k+1}-\rho_{k+1}\,v^k, \qquad \rho_{k+1}:=\sigma_{k+1}/\sigma_k \label{eq:w-reord}
\end{align}
\eqref{eq:w-reord} is exactly $w^{k+1}=2v^{k+1}-v^k$ -- the proposed extrapolation --
whenever $\sigma_k$ is constant ($\rho_{k+1}=1$); the $\rho_{k+1}$ factor is only needed
once $\sigma_k$ genuinely varies, and falls out of the check below rather than being
assumed.

\subsection*{Checking each line}
\textbf{\eqref{eq:u-reord}.} Pure substitution, no algebra: $w^k$'s definition dropped
straight into Part I's $x$-update -- step included, since $\prox_{(1/\tau_k)f_1}$
genuinely depends on it in general,
\[
x^{k+1}=\prox_{(1/\tau_k)f_1}\Big(x^k-\tfrac1{\tau_k}A^\top(\gamma_kr^k-\delta^k)\Big)
=\prox_{(1/\tau_k)f_1}\big(x^k-\tfrac1{\tau_k}A^\top w^k\big)
\ \overset{u=x,\ \theta_k=1/\tau_k}{=}\ \prox_{\theta_kf_1}\big(u^k-\theta_kA^\top w^k\big)
=\eqref{eq:u-reord}.\quad\checkmark
\]
(The $\theta_k=1/\tau_k$ identification makes the \emph{step}, not just the gradient
coefficient, match too -- $\prox_{(1/\tau_k)f_1}=\prox_{\theta_kf_1}$ exactly under that
substitution, no separate argument needed.)

\textbf{\eqref{eq:v-reord}.} Moreau's identity ($\prox_{th}(u):=\arg\min_p h(p)+\tfrac1{2t}\|p-u\|^2$;
verified numerically to machine precision in an earlier pass), for $h:=f_2^*,\ t:=\gamma_k$:
$\prox_{\gamma_kf_2^*}(u)=u-\gamma_k\prox_{f_2/\gamma_k}(u/\gamma_k)$. At
$u:=\gamma_kAx^{k+1}-\delta^k$:
\[
\prox_{\gamma_kf_2^*}\big(\gamma_kAx^{k+1}-\delta^k\big)
=\big(\gamma_kAx^{k+1}-\delta^k\big)-\gamma_k\prox_{f_2/\gamma_k}\big(Ax^{k+1}-\delta^k/\gamma_k\big)
\overset{\text{Part I }y\text{-update}}{=}\gamma_k(Ax^{k+1}-y^{k+1})-\delta^k=-\delta^{k+1},
\]
i.e.\ $v^{k+1}=\prox_{\gamma_kf_2^*}(v^k+\gamma_kAu^{k+1})=\eqref{eq:v-reord}$ with
$\sigma_k=\gamma_k$.\quad$\checkmark$

\textbf{\eqref{eq:w-reord}.} Straight from $w$'s definition at $k+1$:
$w^{k+1}=\gamma_{k+1}r^{k+1}-\delta^{k+1}$. Part I's $\delta$-update gives
$\delta^{k+1}=\delta^k-\gamma_kr^{k+1}$, i.e.\ $r^{k+1}=(\delta^k-\delta^{k+1})/\gamma_k$, so
\[
w^{k+1}=\frac{\gamma_{k+1}}{\gamma_k}(\delta^k-\delta^{k+1})-\delta^{k+1}
=\rho_{k+1}\delta^k-(1+\rho_{k+1})\delta^{k+1}
\overset{\delta=-v}{=}(1+\rho_{k+1})v^{k+1}-\rho_{k+1}v^k=\eqref{eq:w-reord}.\quad\checkmark
\]
(Starting at $k=0$: $w^0:=\gamma_0r^0-\delta^0$ needs no special convention, it's just
the definition evaluated at $k=0$ -- unlike the earlier attempt, no $\delta^{-1}$ or
$\gamma_{-1}$ bookkeeping is needed anywhere.)

\subsection*{An exact equivalence}
\[
\boxed{u^k=x^k,\quad v^k=-\delta^k,\quad w^k=\gamma_kr^k-\delta^k,\quad
\theta_k=\tfrac1{\tau_k},\quad\sigma_k=\gamma_k}
\]
Part I (LADMM, as coded, $\gamma_k,\tau_k$ adaptive) \emph{is}
\eqref{eq:u-reord}--\eqref{eq:w-reord} -- Algorithm 1, reordered -- term for term, index
for index, nothing left over. Algorithm 2's $\theta_{k+1},\sigma_{k+1}$ update formulas
(Part II) now translate directly into $\tau_{k+1}:=1/\theta_{k+1}$,
$\gamma_{k+1}:=\sigma_{k+1}$, \emph{and} into \eqref{eq:w-reord}'s $\rho_{k+1}$ -- the
one thing to carry along that their own paper doesn't need, since their un-reordered
Algorithm 1 never has a second stored value of the dual variable to combine.

\subsection*{Where did $y^k$ go?}
Not lost -- Moreau's identity above only \emph{needed} the difference
$\gamma_k(Ax^{k+1}-y^{k+1})-\delta^k$, but it was computed by literally evaluating
$\prox_{f_2/\gamma_k}(\cdot)$ along the way, and that inner evaluation \emph{is}
$y^{k+1}$ (Part I's own $y$-update, verbatim). Reading it back off, using
$\prox_{\gamma_kf_2^*}(u)=u-\gamma_ky^{k+1}=-\delta^{k+1}$ with $u=\gamma_kAx^{k+1}-\delta^k$:
\[
\gamma_k y^{k+1} = u+\delta^{k+1} = \big(\gamma_kAx^{k+1}-\delta^k\big) + \delta^{k+1}
\quad\Longrightarrow\quad
y^{k+1} = Ax^{k+1} + \frac{\delta^{k+1}-\delta^k}{\gamma_k},
\]
and substituting $\delta=-v$:
\begin{equation}
\boxed{y^{k+1} = Au^{k+1} + \frac{v^k-v^{k+1}}{\gamma_k} = Au^{k+1}+\frac{v^k-v^{k+1}}{\sigma_k}} \label{eq:yrelation}
\end{equation}
(shift the index down by one for $y^k$ in terms of iteration-$(k-1)$ quantities:
$y^k=Au^k+(v^{k-1}-v^k)/\sigma_{k-1}$, valid $k\ge1$). Sanity check against Part II's own
$d_{k+1}$ residual (\S\,Part II): $d_{k+1}=|\tfrac1{\sigma_k}(v^k-v^{k+1})-A(u^k-u^{k+1})|$
-- visibly the same $\tfrac1{\sigma_k}(v^k-v^{k+1})$ term as in \eqref{eq:yrelation}, but
paired with $A(u^k-u^{k+1})$ there instead of $Au^{k+1}$ here; they are related but not
the same quantity, worth not conflating.

So: $x^k=u^k$, $\delta^k=-v^k$, $w^k=\gamma_kr^k-\delta^k$, $\tau_k=1/\theta_k$,
$\gamma_k=\sigma_k$, and $y^k$ via \eqref{eq:yrelation} --
every quantity in Part I now has an exact counterpart in Algorithm 1, reordered, and
vice versa. What's still open is not the equivalence itself but whether Algorithm 2's
$p_k,d_k$ residuals and step-size update -- derived and proved for their own
(un-reordered) Algorithm 1 -- remain valid (still drive $\theta_k,\sigma_k$ toward a
convergent, stabilizing choice) once read against \eqref{eq:u-reord}--\eqref{eq:w-reord}
instead. Not claimed either way here -- next thing to check.

\section*{Part IV -- can the extrapolation coefficient itself be fixed?}
Part III's equivalence is exact, but its extrapolation line needs a genuinely adaptive
coefficient, $w^{k+1}=(1+\rho_{k+1})v^{k+1}-\rho_{k+1}v^k$ with
$\rho_{k+1}=\gamma_{k+1}/\gamma_k$ -- not Algorithm 2's own literal, fixed $(2,-1)$.
Question: is there a reindexing that restores the fixed form exactly, for genuinely
adaptive $\gamma_k,\tau_k$? Answer, as far as checked so far: yes, exactly one
construction does it, and it comes at a real cost, not for free.

\subsection*{The fix: lag the $x$-update's own step by one iteration}
Leave the $y$- and $\delta$-updates exactly as coded, using $\gamma_k$ (governing the
$k\to k{+}1$ transition, same as always). Change only the $x$-update, to use the
\emph{previous} iteration's step sizes -- \emph{both} occurrences, the prox's own step
and the gradient coefficient, since (in general) they're tied together and shifting one
without the other would silently be a different, uncontrolled algorithm:
\[
x^{k+1} := \prox_{(1/\tau_{k-1})f_1}\Big(x^k - \tfrac1{\tau_{k-1}}A^\top\big(\gamma_{k-1}r^k-\delta^k\big)\Big)
\qquad(\text{was }\gamma_k,\tau_k\text{ -- see Part I}).
\]
Define $w^k:=\gamma_{k-1}r^k-\delta^k$, matching this new bracket. The \emph{unshifted}
$\delta$-update already gives $\delta^k=\delta^{k-1}-\gamma_{k-1}r^k$, i.e.
$\gamma_{k-1}r^k=\delta^{k-1}-\delta^k$ directly -- no ratio factor, because both sides
now refer to the same $\gamma_{k-1}$. So
\[
w^k=(\delta^{k-1}-\delta^k)-\delta^k=\delta^{k-1}-2\delta^k
\overset{\delta=-v}{=}2v^k-v^{k-1},
\]
Algorithm 2's literal, unmodified extrapolation -- for \emph{any} step-size sequence,
adaptive or not, no special case needed.

\subsection*{The full correspondence, for this lagged version}
Boundary convention first, spelled out rather than left implicit: $w^0$ needs
$\gamma_{-1},\tau_{-1}$ (there is no ``iteration $-1$''), so set
$\gamma_{-1}:=\gamma_0,\ \tau_{-1}:=\tau_0$ -- the natural ``no history yet, first step
is unshifted'' choice, giving $w^0=\gamma_0r^0-\delta^0$ (Part III's own $w^0$) and
$v^{-1}:=v^0$, so $w^0=2v^0-v^{-1}=v^0$: no real extrapolation on the very first step,
same as Part III.

With $u^k:=x^k,\ v^k:=-\delta^k$ as before, matching this lagged scheme against
``Algorithm 1, reordered'' (\eqref{eq:u-reord}--\eqref{eq:w-reord}'s template, now with
the extrapolation fixed at $(2,-1)$) gives
\[
\boxed{u^k=x^k,\quad v^k=-\delta^k,\quad w^k=\gamma_{k-1}r^k-\delta^k,\quad
\theta_k=\frac1{\tau_{k-1}},\quad \sigma_k=\gamma_k}
\]
Look closely at the last two: $\sigma_k$ uses $\gamma_k$ directly (unshifted -- the
$y,\delta$-update was never touched), but $\theta_k$ uses $\tau_{k-1}$ (shifted). That
index mismatch \emph{is} \S~``What this costs'' point 1, made concrete: Algorithm 2's
own rule produces $\theta_{k+1},\sigma_{k+1}$ as a \emph{matched pair} from the same
$p_{k+1},d_{k+1}$, meant to govern the same transition -- here $\sigma_{k+1}=\gamma_{k+1}$
governs producing $\delta^{k+2}$ from $\delta^{k+1}$, while $\theta_{k+1}=1/\tau_k$
doesn't get used until producing $x^{k+3}$ from $x^{k+2}$ -- one physical LADMM
iteration later than its partner. Algorithm 2 was not written with that mismatch in
mind; nothing here says it still behaves sensibly under it.

$y^k$'s relation to $u,v$ is untouched by any of this -- it only ever depended on the
(unshifted) $y,\delta$-update, so Part III's \eqref{eq:yrelation} carries over verbatim:
$y^{k+1}=Au^{k+1}+(v^k-v^{k+1})/\sigma_k$.

\subsection*{Primal/dual residuals, adaptive step, lagged scheme}
Derived directly from this problem's own feasibility/stationarity conditions -- not
borrowed from \code{ladmm.py}'s or Goldstein--Li--Yuan's own residual constructions.
The Lagrangian sign convention that reproduces the already-established
$-\delta^\star\in\partial f_2(y^\star)$ is
$L_\gamma=f_1(x)+f_2(y)-\langle\delta,Ax-y\rangle+\tfrac\gamma2\|Ax-y\|^2$, giving
\[
\textbf{primal feasibility: } Ax^\star=y^\star, \qquad
\textbf{dual feasibility: } A^\top\delta^\star\in\partial f_1(x^\star),\ \
-\delta^\star\in\partial f_2(y^\star).
\]
\textbf{Primal residual:} trivially $D^k:=Ax^k-y^k$ -- \code{ladmm.py}'s own $D$ (and,
by this standard terminology, genuinely the \emph{primal} one, reinforcing the earlier
naming note).

\textbf{Dual residual, $y$-side: exactly zero, always.} $y^{k+1}$ solves its
subproblem exactly, and $\delta^{k+1}$ is defined so that its own optimality condition
\emph{is} $-\delta^{k+1}\in\partial f_2(y^{k+1})$ -- not approximately, exactly, every
iteration. All the real slack is on the $x$-side, where $x$ is only linearized.

\textbf{Dual residual, $x$-side.} The lagged $x$-update's own optimality condition,
matching \code{ladmm.py}'s actual gradient bracket exactly:
\[
\tau_{k-1}(x^k-x^{k+1}) - \gamma_{k-1}A^\top D^k + A^\top\delta^k \in \partial f_1(x^{k+1}).
\]
Subtracting the true target $A^\top\delta^{k+1}\in\partial f_1(x^{k+1})$, and
$\delta^k-\delta^{k+1}=\gamma_kD^{k+1}$:
\begin{equation}
\boxed{P^{k+1}_{\mathrm{lagged}} = \tau_{k-1}(x^k-x^{k+1}) + A^\top\big(\gamma_kD^{k+1}-\gamma_{k-1}D^k\big)} \label{eq:P-kkt-lagged}
\end{equation}
Directly traceable to this problem's own feasibility conditions, line by line, using
this scheme's own updates -- nothing borrowed. (Reduces, when $\gamma,\tau$ are
constant, to $\tau(x^{k-1}-x^k)+\gamma A^\top(D^k-D^{k-1})$ after reindexing -- close to
but not identical to \code{ladmm.py}'s own coded $P^k=\tau(x^{k-1}-x^k)+\gamma
A^\top(D^{k-1}-D^k)$, opposite sign on the second term. Not a contradiction: they're
different constructions from different proofs, both vanishing at a true fixed point;
\eqref{eq:P-kkt-lagged} is the one with a fully transparent derivation traceable to
this problem's own KKT conditions.)

\subsection*{Why the other direction doesn't work}
The mirror idea -- leave the $x$-update alone, shift the $y,\delta$-update
\emph{forward} to use $\gamma_{k+1}$ instead of $\gamma_k$ -- gives the same clean
algebra, but isn't usable: $\gamma_{k+1}$, chosen (by Algorithm 2's own rule) from
residuals $p_{k+1},d_{k+1}$, needs $v^{k+1}=-\delta^{k+1}$ -- exactly the quantity that
shifted $\delta$-update is trying to produce. Circular. The backward shift above has no
such problem: $\gamma_{k-1}$ was already fixed, from residuals computed a full
iteration earlier, well before the $x$-update at step $k$ runs.

\subsection*{What this costs}
Two things, both real, neither a labeling issue:
\begin{enumerate}
\item \textbf{An extra lag.} In Algorithm 2's own convention, $\theta_j$ (decided from
iteration-$j$ residuals) is used immediately, to produce $u^{j+1}$. Here, $\tau_j$
(decided the same way, from the same residuals) isn't used until producing $x^{j+2}$ --
one iteration later than Algorithm 2's own rule assumes for its own step size. This is
a structural change to the algorithm, not a renaming: $\tau$'s feedback loop is now one
step slower than $\gamma$'s.
\item \textbf{Still not literally their Algorithm 1.} This shifted scheme is exactly
``Algorithm 1, reordered'' (Part III) with a fixed extrapolation coefficient --
but ``reordered'' extrapolates $v$ into the $u$-update, while Goldstein--Li--Yuan's own,
published Algorithm 1 extrapolates $u$ into the $v$-update. These are different
recursions in general (not equal given the same $u^k,v^k$, except at a fixed point), so
their convergence proof -- and Algorithm 2's proof that its adaptive rule stays
convergent -- was written for \emph{their} ordering, not this one. Nothing here shows
the shifted scheme converges under Algorithm 2's step-size rule; it only shows that,
\emph{if} it converges, its extrapolation term has exactly the form they wrote down.
\end{enumerate}

\subsection*{Where this leaves things}
Of the two ways to try to reconcile adaptive $\gamma_k,\tau_k$ with Algorithm 2's
literal, fixed extrapolation, only one is even well-defined: lagging $\tau$ one
iteration behind $\gamma$ is causal; shifting the other way needs a residual before the
quantity it depends on exists. So, as far as checked here: \textbf{this one-line,
backward index shift on the $x$-update is the only way found so far to make
adaptive-step LADMM and adaptive-step PDHG use the exact same, unmodified extrapolation
formula.} It is not, by itself, a proof that the resulting scheme converges -- that's
still open, and separate from Part III's exact (but coefficient-adaptive) equivalence,
which needs no such shift and no such open question.

\section*{Part V -- what does the adaptive-LADMM literature actually do?}
Checked (web search, not exhaustive) rather than assumed. Two representative papers
that adapt a linearized-ADMM-style penalty parameter:
\begin{itemize}
\item Lin, Liu \& Su, ``Linearized Alternating Direction Method with Adaptive Penalty
for Low-Rank Representation,'' NeurIPS 2011 (arXiv:1109.0367) -- LADMAP. Linearizes one
subproblem exactly the way Part I does; adapts the penalty from a feasibility/residual
condition. \textbf{No extrapolation or momentum term anywhere} -- no $2x^{k+1}-x^k$, no
lagged combination of past dual iterates. Just Part I's update, run with $\gamma_k$
adapted, unmodified otherwise.
\item Xu, Figueiredo \& Goldstein, ``Adaptive ADMM with Spectral Penalty Parameter
Selection,'' AISTATS 2017 (arXiv:1605.07246) -- penalty adapted each iteration via a
Barzilai--Borwein spectral estimate from the primal/dual residuals. Every description
found of it (full text wasn't extractable to triple-check) is again a residual-driven
penalty update on the standard ADMM/LADMM subproblems, with no mention of an
extrapolation step.
\end{itemize}
Neither paper frames its adaptivity in terms of PDHG or Chambolle--Pock at all -- the
question of whether their adapted iteration still matches \emph{some} adaptive PDHG
scheme isn't asked.

\paragraph{So, following the logic of Parts III--IV:} both are exactly the ``no lag''
LADMM of Part I, with $\gamma_k,\tau_k$ adapted directly, no shift. By Part III, that
recursion \emph{is} Algorithm~1, reordered, with the \emph{generalized}
$(1+\rho_{k+1},-\rho_{k+1})$ extrapolation -- not Goldstein--Li--Yuan's own literal,
fixed $(2,-1)$ Algorithm~2. So, as best determined here: \textbf{the adaptive-LADMM
scheme actually used in this literature is a genuinely different algorithm from
adaptive PDHG (Algorithm 2) whenever the step size actually changes between
iterations, and coincides with it only in the constant-step case} ($\rho_k\equiv1$,
Part III's original, pre-adaptive equivalence). That specific comparison -- LADMM's own
common adaptive-penalty practice against Goldstein--Li--Yuan's adaptive PDHG,
specifically -- wasn't found stated anywhere in the two papers surveyed.

\paragraph{Closest related idea found, for honesty's sake.} Malitsky's golden-ratio /
reflected-gradient line of work (e.g.\ ``Golden Ratio Algorithms for Variational
Inequalities,'' 2018/2020) does carefully track exactly this kind of issue in a related
setting: their reflected/extrapolated schemes are proven equivalent to the relevant
primal-dual method \emph{under constant step sizes}, and their adaptive-step variants
build the step size and the extrapolation from the same local history (consecutive
iterates) rather than treating them independently -- structurally the same flavor of
fix as Part IV's, though for a different family of methods (variational-inequality
reflected gradient, not linearized ADMM specifically) and not verified here to reduce
to Part III's or IV's formulas exactly. So the phenomenon itself -- adaptive steps
force the extrapolation/reflection term to track step-size history -- isn't
unprecedented in the broader primal-dual literature; what wasn't found is this specific
comparison stated for LADMM's own common adaptive-penalty practice against PDHG.

\section*{Part VI -- dualizing the dual: a second route to Algorithm 1}
A different idea for resolving Part III's ordering mismatch: instead of forming the
saddle point from the \emph{primal} problem (Part II/III's route), form it from the
\emph{dual} problem instead, and see if the resulting Algorithm 1 instance already
matches Part I without any hand-built reordering. Worked through in full below --
genuinely different from Part III's route, not assumed to land in the same place.
New symbols throughout ($\tilde u,\tilde v,\tilde\theta,\tilde\sigma$), so nothing here
is confused with Part II/III's own $u,v,\theta,\sigma$.

\subsection*{The dual problem, and its own saddle point}
Primal: $\min_u f_1(u)+f_2(Au)$, $u=x$ (Part II/III). Its Fenchel dual (general
formula: $\min_x f(x)+g(Kx)$ has dual $\min_p f^*(-K^\top p)+g^*(p)$, derived earlier
this session) is
\[
\min_{\tilde u}\ f_2^*(\tilde u) + f_1^*(-A^\top\tilde u).
\]
Legendre-transform \emph{this} problem's own nonsmooth-composed piece
($f_1^*(-A^\top\tilde u)$, using $(f_1^*)^*=f_1$) to get a saddle point, matching
Goldstein's template ($\min_X\max_V F(X)+\langle V,KX\rangle-G(V)$) with
$X=\tilde u,\ F=f_2^*,\ K=-A^\top,\ G=f_1$:
\[
\min_{\tilde u}\max_{\tilde v}\ f_2^*(\tilde u) + \langle\tilde v,-A^\top\tilde u\rangle - f_1(\tilde v).
\]

\subsection*{Side by side}
\begin{center}
\begin{tabular}{@{}lll@{}}
\toprule
& Part II (primal-dualized) & Part VI (dual-redualized)\\
\midrule
Saddle point & $\min_u\max_v f_1(u)+\langle v,Au\rangle-f_2^*(v)$ & $\min_{\tilde u}\max_{\tilde v} f_2^*(\tilde u)+\langle\tilde v,-A^\top\tilde u\rangle-f_1(\tilde v)$\\
$X$ (no extrap., updates first) & $u$, role $F=f_1$ & $\tilde u$, role $F=f_2^*$\\
$V$ (uses extrap.\ fresh $X$) & $v$, role $G=f_2^*$ & $\tilde v$, role $G=f_1$\\
Operator & $A$ & $-A^\top$\\
\bottomrule
\end{tabular}
\end{center}
Genuinely different from Part II's instance -- which of $f_1,f_2^*$ plays $F$ vs.\ $G$
is swapped, and the operator is $-A^\top$, not $A$. Algorithm 1's template, applied here:
\begin{align}
\tilde u^{k+1} &= \prox_{\tilde\theta_k f_2^*}\big(\tilde u^k+\tilde\theta_kA\tilde v^k\big) \label{eq:ut}\\
\tilde v^{k+1} &= \prox_{\tilde\sigma_k f_1}\big(\tilde v^k-\tilde\sigma_kA^\top(2\tilde u^{k+1}-\tilde u^k)\big) \label{eq:vt}
\end{align}

\subsection*{Checking against Part I}
\textbf{\eqref{eq:ut}, the $y/\delta$-side.} Try $\tilde v^k:=x^{k+1}$ (forced -- see
below), $\tilde u^k:=-\delta^k$: \eqref{eq:ut} reads
$-\delta^{k+1}=\prox_{\tilde\theta_kf_2^*}(-\delta^k+\tilde\theta_kAx^{k+1})$, which
matches Part III's own $y/\delta$-update (via the same Moreau step, unchanged) exactly
with $\tilde\theta_k=\gamma_k$ -- \textbf{no lag, no $\rho$-correction, clean.} This is
a real improvement over Part IV's construction on this side.

(Why $\tilde v^k=x^{k+1}$, not $x^k$: Algorithm 1's own order updates $X$ first, using
whatever $V$ is already stored -- so \eqref{eq:ut} can only ever see a \emph{stored}
$\tilde v$, never a freshly-computed one. Part I's $y/\delta$-update needs the
\emph{fresh} $x^{k+1}$ (computed by the $x$-update earlier in the same LADMM
iteration), so the only consistent identification is $\tilde v$ running one index
ahead of $x$ -- forced by matching terms, not a free choice.)

\textbf{\eqref{eq:vt}, the $x$-side.} With that same identification,
$\tilde v^{k+1}=x^{k+2}$, and \eqref{eq:vt} becomes
$x^{k+2}=\prox_{\tilde\sigma_kf_1}\big(x^{k+1}+\tilde\sigma_kA^\top(2\delta^{k+1}-\delta^k)\big)$.
Matching against Part III's general (adaptive) $x$-update at the corresponding index
forces $\tilde\sigma_k=\rho_{k+1}/\tau_{k+1}$ \emph{and} $2\tilde\sigma_k=(1+\rho_{k+1})/\tau_{k+1}$
simultaneously -- solvable only when $\rho_{k+1}=1$, i.e.\ $\gamma_{k+1}=\gamma_k$. Same
verdict as Part IV, reached a different way: \eqref{eq:vt}'s extrapolation is fixed at
$(2,-1)$ (built into Algorithm 1's template, unmodified, on this side now), and fixed
coefficients only reproduce fully adaptive-step LADMM when the step happens to be
constant over that pair of iterations.

Worse, the fix isn't obviously free here either: making \eqref{eq:vt} exact for
adaptive steps would need $\tilde\sigma_k=\rho_{k+1}/\tau_{k+1}$, i.e.\ knowing
$\gamma_{k+1}$ \emph{before} computing $\tilde v^{k+1}$ -- and $\gamma_{k+1}$ would
ordinarily be chosen from residuals of a transition this construction's own step order
hasn't reached yet at that point. This has the same shape as Part IV's rejected
``shift the other direction'' argument, though I haven't nailed down with full
certainty whether it's actually circular here or merely looks that way by analogy --
flagged as open, not asserted.

\subsection*{Verdict}
Not the same algorithm as Part II's, and not a way to escape Part III/IV's tension
either -- it relocates it. Genuinely useful outcome, though: it cleanly fixes \emph{half}
the mismatch (the $y/\delta$-side needs no lag or correction at all, an improvement
over every earlier construction), leaving only the $x$-side facing the same
fixed-vs-adaptive-extrapolation issue Part IV already found. Reduces exactly to Part
IV's constant-step case when $\gamma_k\equiv\gamma$. Worth having on record as a second,
independent route to the same underlying obstacle -- not a resolution of it.
\end{document}
